% !TEX root = ../../summer_2026_introduction.tex
  \subsection{\textbf{満点を取らなければいけない？}}
    結論から言おう。\textbf{入試において満点を取らなければいけないわけがない。}

    何故ならば、\textbf{入試は合計点勝負}であって、特定の教科において下限が設定されているわけではないからである。
    極端に言うと、1教科0点を取ったからといって必ずしも不合格になるという道理はないのである。

  \subsection{\textbf{高校進学は人生の通過点！}}
    みんなは、まだまだ若いけど、今この時期が人生の大半を左右するかもしれない時期にあることは理解しておいた方がいい。

    長い人生、今よりキツい時が必ず来るけど、今を乗り越えると未来の自分に感謝されると思うよ!!

  \subsection{\textbf{ちょっと小話をしよう}}
  「分かる」とはそもそもどういうことでしょうか。
  よく「分かる」とは「分ける」ことだと言われます。みなさんの頭の中に沢山のラベルがつけられた「箱」があるとイメージしてみてください。みなさんは郵便局員のように、入ってくる情報に対して，これはこっちの箱，これはこっちの箱と仕分けをしていきます。
  新しい情報を頭の中のしかるべき「箱」に入れることが「分かる」ことであり、どの箱にいれるべきかを的確にガイドしてくれる説明がいわゆる「分かりやすい説明」というわけです。
  しかし。もしやってきた情報が、みなさんにとって全く新しい場意であったらどうでしょう。それはみなさんの持っているどの箱にも分類できないようなものです。「分けられない」のですから『分からない」。その「分からない」状態でいることに私達は我慢できなくなってしまうのですね。そんなとき、その概念を手近な箱に無理やり押し込んでくれる説明があると，それを「分かりやすい」と感じてしまいがちです。でもそれはみなさんを安心させるためのその場しのぎの説明。確かに分かったつもりにはなれるでしょうが、誤った箱に入れられた情報は二度と顧みられることはなく、本質的
  「分かりやすさ」のとても危険な落とし穴なのです。

  においてとても大切な姿勢は何か。それは「分からないを分からないものとして一旦受け入れるということ」だと僕は思うのです。言ってみれば「分からない」というラベルのついた「箱」を作り、行き場のない情報はひとまずそこに放り込んでおくということです。
  もちろんそれは少し居心地が悪いことかもしれません。でもその違和感をうまくいならしながら、ときどきは箱の中をチラチラと眺めながら先に進んでいくと、不思議なもので「分からない」の箱にいれておいたいくつかの情報が結びつき、きちんと意味を持ちはじめる日がくるのです。頭の中にその概念をおさめる新しい「箱」が立ち上がる。これが真の「エウレカ！（分かったぞ！）」の瞬間です。
  「分からない」は学びの障害物ではありません。むしろ学びにおいてなくてはならない重要な通過点なのです。
  

  \newpage
